\chapter{Architecture}
\label{ch:architecture}

\section{Design Philosophy}

\newthought{OutOf5K follows} several core architectural principles that guide its design:

\begin{description}
    \item[Local-First] All demo processing happens locally on the user's machine, eliminating server costs and preserving user privacy.
    
    \item[Modular Architecture] The system is composed of loosely coupled components that communicate through well-defined interfaces.
    
    \item[Progressive Enhancement] Start with essential features and progressively add advanced capabilities.
    
    \item[Performance-Conscious] CS2 demo files can be large (100MB+). The architecture handles this through background processing and incremental parsing.
\end{description}

\section{System Context}

\begin{figure*}[h]
    \centering
    \includegraphics[width=0.9\textwidth]{figures/c4/1-context}
    \caption{C4 Context Diagram showing OutOf5K in relation to external systems}
    \label{fig:c4-context}
\end{figure*}

\marginnote{The C4 model provides a hierarchical way to describe software architecture at different levels of abstraction.}

\outof{} operates as a standalone desktop application that interacts with several external systems:

\begin{itemize}
    \item \textbf{Steam Platform}: User authentication and match history
    \item \textbf{HLTV}: Professional player statistics and match data
    \item \textbf{Faceit/ESEA}: Competitive match data and rankings
    \item \textbf{CS2 Game}: Demo file source
\end{itemize}

\section{Container Architecture}

The application is composed of four main containers, as shown in Figure~\ref{fig:c4-container}.

\begin{figure*}[h]
    \centering
    \includegraphics[width=0.9\textwidth]{figures/c4/2-container}
    \caption{C4 Container Diagram showing the major components}
    \label{fig:c4-container}
\end{figure*}

\subsection{Frontend (Svelte + TypeScript)}

The user interface layer built with Svelte 5 and TypeScript. Key responsibilities:

\begin{itemize}
    \item Rendering UI components
    \item User interaction handling
    \item 2D map visualization via Pixi.js
    \item Statistics charts and graphs
    \item State management via Svelte stores
\end{itemize}

\subsection{Tauri Backend (Rust)}

The desktop application shell and system integration layer:

\begin{itemize}
    \item Native window management
    \item File system access for demo files
    \item SQLite database operations
    \item Python subprocess management
    \item IPC command handlers
\end{itemize}

\subsection{Python Service}

A subprocess handling computationally intensive demo parsing:

\begin{itemize}
    \item CS2 demo file parsing via demoparser2
    \item Data extraction and transformation
    \item Heavy computational analysis
    \item Future: ML model inference
\end{itemize}

\subsection{SQLite Database}

Local persistent storage:

\begin{itemize}
    \item Parsed demo metadata and events
    \item Player statistics and profiles
    \item User preferences and settings
    \item Cached external API data
\end{itemize}

\section{Component Details}

\subsection{Frontend Components}

\begin{figure}[h]
    \centering
    \includegraphics[width=\linewidth]{figures/c4/3-component-frontend}
    \caption{Frontend component diagram}
    \label{fig:c4-frontend}
\end{figure}

The frontend is organized into several key components:

\begin{description}
    \item[Router] Handles page navigation using SvelteKit's routing
    \item[Dashboard] Main landing page with activity feed
    \item[Demo Library] Lists all imported demos with filtering
    \item[Demo Viewer] Main analysis view with timeline and map
    \item[Map Canvas] Pixi.js-based 2D rendering
    \item[Stores] Svelte stores for state management
    \item[API Layer] Typed wrappers around Tauri IPC
\end{description}

\subsection{Rust Backend Components}

\begin{figure}[h]
    \centering
    \includegraphics[width=\linewidth]{figures/c4/3-component-rust}
    \caption{Rust backend component diagram}
    \label{fig:c4-rust}
\end{figure}

Key backend components:

\begin{description}
    \item[Commands] Tauri IPC entry points
    \item[Services] Business logic orchestration
    \item[Repositories] Data access layer
    \item[Python Bridge] Subprocess communication
    \item[API Clients] External service integration
\end{description}

\subsection{Python Service Components}

\begin{figure}[h]
    \centering
    \includegraphics[width=\linewidth]{figures/c4/3-component-python}
    \caption{Python service component diagram}
    \label{fig:c4-python}
\end{figure}

The Python service includes:

\begin{description}
    \item[Main Loop] Command dispatch and response handling
    \item[Demo Parser] demoparser2 integration
    \item[Event Processor] Game event classification
    \item[Stats Calculator] Performance metrics
    \item[Highlight Detector] Notable moment detection
\end{description}

\section{Technology Rationale}

\begin{table}[h]
    \centering
    \begin{tabular}{@{}ll@{}}
        \toprule
        \textbf{Decision} & \textbf{Rationale} \\
        \midrule
        Tauri & Small bundles, native performance \\
        Svelte & Minimal boilerplate, reactive \\
        Python subprocess & Best demo parser available \\
        SQLite & Zero-config, file-based \\
        Pixi.js & High-performance 2D canvas \\
        \bottomrule
    \end{tabular}
    \caption{Key technology decisions}
\end{table}

Detailed rationale for each decision is documented in the Architecture Decision Records (Chapter~\ref{ch:decisions}).
